% GENERATED FILE. Edit canonical lessons, then run npm run generate:books.
% canonical-source-hash: fnv1a64:a589a32c16689d77
% canonical-lessons: KA-C346-sphota, KA-C346-tattarisu, KA-C346-kampisu, KA-C346-pralaya, KA-C346-gayalu

\chapter{An Explosion and a Shock}
\label{ch:ka-an-explosion-and-a-shock}
\hlchaptermodality{346}

\begin{chapteropening}
\textbf{By the end of this chapter:} \emph{I can say an explosion, to reel, to tremble, a deluge and an injured person.}

Complete the last lesson of chapter 346: I can say an explosion, to reel, to tremble, a deluge and an injured person.
\end{chapteropening}

\section[\texorpdfstring{\kn{ಸ್ಫೋಟ} (sphōṭa)}{ಸ್ಫೋಟ (sphōṭa)}]{\kn{ಸ್ಫೋಟ} (sphōṭa) — an explosion; a blast}

\label{lesson:KA-C346-sphota}



\begin{quote}
Before the new one: say the Kannada for a spy, then the Kannada for to sneak in.
\end{quote}

\subsection*{You'll want to know: \kn{ಸ್ಫೋಟ}}
\textbf{\kn{ಸ್ಫೋಟ}} — \emph{sphōṭa} — \textquotedblleft{}an explosion; a blast\textquotedblright{}.

\begin{quote}
\kn{ಕಾರ್ಖಾನೆಯಲ್ಲಿ} \kn{ದೊಡ್ಡ} \kn{ಸ್ಫೋಟ} \kn{ಆಯಿತು}.
\end{quote}

\begin{quote}
\emph{kārkhāneyalli doḍḍa sphōṭa āyitu.}
\end{quote}

\begin{quote}
There was a big explosion at the factory.
\end{quote}

Plural: \kn{ಸ್ಫೋಟಗಳು}, \emph{sphōṭagaḷu}.

\kn{ಸ್ಫೋಟಿಸು} is to explode or to blow up.

\begin{cousinweb}
\textbf{\kn{ಸ್ಫೋಟ}} (\emph{sphōṭa}) is borrowed from Sanskrit \emph{sphoṭa}, \textquotedblleft{}bursting, opening, expansion\textquotedblright{} (Monier-Williams).
\end{cousinweb}

\begin{grammarlens}[title={What you've built}]
One more word for telling what happened.
\end{grammarlens}

\subsection*{Your turn}
\begin{itemize}
\raggedright
  \item \emph{Say it:} \emph{sphōṭa}
  \item \emph{Say it:} \emph{sphōṭa}, once more
  \item \emph{Say it:} \emph{kārkhāneyalli doḍḍa sphōṭa āyitu.} — There was a big explosion at the factory.
  \item \emph{Recall:} say the Kannada for a spy, then the Kannada for to sneak in, then say \emph{sphōṭa} again
\end{itemize}

\begin{tcolorbox}[breakable,colback=teal!4,colframe=teal!35!black,title={Before you move on}]
What does \kn{ಸ್ಫೋಟ} mean? (\textquotedblleft{}An explosion; a blast\textquotedblright{}.) Say it once more.
\end{tcolorbox}
\section[\texorpdfstring{\kn{ತತ್ತರಿಸು} (tattarisu)}{ತತ್ತರಿಸು (tattarisu)}]{\kn{ತತ್ತರಿಸು} (tattarisu) — to reel; to be shaken}

\label{lesson:KA-C346-tattarisu}



\begin{quote}
Before the new one: say the Kannada for to sneak in, then the Kannada for an explosion.
\end{quote}

\subsection*{You'll want to know: \kn{ತತ್ತರಿಸು}}
\textbf{\kn{ತತ್ತರಿಸು}} — \emph{tattarisu} — \textquotedblleft{}to reel; to be shaken\textquotedblright{}.

\begin{quote}
\kn{ಆ} \kn{ಸುದ್ದಿಯಿಂದ} \kn{ಇಡೀ} \kn{ಊರು} \kn{ತತ್ತರಿಸಿತು}.
\end{quote}

\begin{quote}
\emph{ā suddiyinda iḍī ūru tattarisitu.}
\end{quote}

\begin{quote}
The whole town reeled at that news.
\end{quote}

Past (it …): \kn{ತತ್ತರಿಸಿತು}, \emph{tattarisitu}.

A place or a person shaken by a blow or a loss.

\begin{cousinweb}
\textbf{\kn{ತತ್ತರಿಸು}} (\emph{tattarisu}) is native Dravidian: DEDR 3061 lists Kannada \emph{tattaṟa}, \emph{tattaṟi}, \textquotedblleft{}perplexity, confusion\textquotedblright{}, beside Tamil \emph{tattaram}, \textquotedblleft{}tremulousness, quivering, flurry\textquotedblright{}, here with the verb ending -\emph{isu}.
\end{cousinweb}

\begin{grammarlens}[title={What you've built}]
One more word for telling what happened.
\end{grammarlens}

\subsection*{Your turn}
\begin{itemize}
\raggedright
  \item \emph{Say it:} \emph{tattarisu}
  \item \emph{Say it:} \emph{tattarisu}, once more
  \item \emph{Say it:} \emph{ā suddiyinda iḍī ūru tattarisitu.} — The whole town reeled at that news.
  \item \emph{Recall:} say the Kannada for to sneak in, then the Kannada for an explosion, then say \emph{tattarisu} again
\end{itemize}

\begin{tcolorbox}[breakable,colback=teal!4,colframe=teal!35!black,title={Before you move on}]
What does \kn{ತತ್ತರಿಸು} mean? (\textquotedblleft{}To reel; to be shaken\textquotedblright{}.) Say it once more.
\end{tcolorbox}
\section[\texorpdfstring{\kn{ಕಂಪಿಸು} (kampisu)}{ಕಂಪಿಸು (kampisu)}]{\kn{ಕಂಪಿಸು} (kampisu) — to tremble; to shake}

\label{lesson:KA-C346-kampisu}



\begin{quote}
Before the new one: say the Kannada for an explosion, then the Kannada for to reel.
\end{quote}

\subsection*{You'll want to know: \kn{ಕಂಪಿಸು}}
\textbf{\kn{ಕಂಪಿಸು}} — \emph{kampisu} — \textquotedblleft{}to tremble; to shake\textquotedblright{}.

\begin{quote}
\kn{ಸ್ಫೋಟದಿಂದ} \kn{ಹತ್ತಿರದ} \kn{ಮನೆಗಳು} \kn{ಕಂಪಿಸಿದವು}.
\end{quote}

\begin{quote}
\emph{sphōṭadinda hattirada manegaḷu kampisidavu.}
\end{quote}

\begin{quote}
The houses nearby shook from the blast.
\end{quote}

Past (it …): \kn{ಕಂಪಿಸಿತು}, \emph{kampisitu}.

The ground, a building or a voice trembles: \kn{ಭೂಮಿ} \kn{ಕಂಪಿಸಿತು}.

\begin{cousinweb}
\textbf{\kn{ಕಂಪಿಸು}} (\emph{kampisu}) is built on Sanskrit \emph{kampa}, \textquotedblleft{}trembling, a tremor\textquotedblright{} (CDIAL 2765), with the Kannada verb ending -\emph{isu}.
\end{cousinweb}

\begin{grammarlens}[title={What you've built}]
One more word for telling what happened.
\end{grammarlens}

\subsection*{Your turn}
\begin{itemize}
\raggedright
  \item \emph{Say it:} \emph{kampisu}
  \item \emph{Say it:} \emph{kampisu}, once more
  \item \emph{Say it:} \emph{sphōṭadinda hattirada manegaḷu kampisidavu.} — The houses nearby shook from the blast.
  \item \emph{Recall:} say the Kannada for an explosion, then the Kannada for to reel, then say \emph{kampisu} again
\end{itemize}

\begin{tcolorbox}[breakable,colback=teal!4,colframe=teal!35!black,title={Before you move on}]
What does \kn{ಕಂಪಿಸು} mean? (\textquotedblleft{}To tremble; to shake\textquotedblright{}.) Say it once more.
\end{tcolorbox}
\section[\texorpdfstring{\kn{ಪ್ರಳಯ} (praḷaya)}{ಪ್ರಳಯ (praḷaya)}]{\kn{ಪ್ರಳಯ} (praḷaya) — a deluge; a catastrophe}

\label{lesson:KA-C346-pralaya}



\begin{quote}
Before the new one: say the Kannada for to reel, then the Kannada for to tremble.
\end{quote}

\subsection*{You'll want to know: \kn{ಪ್ರಳಯ}}
\textbf{\kn{ಪ್ರಳಯ}} — \emph{praḷaya} — \textquotedblleft{}a deluge; a catastrophe\textquotedblright{}.

\begin{quote}
\kn{ಆ} \kn{ರಾತ್ರಿ} \kn{ಪ್ರಳಯ} \kn{ಬಂದ} \kn{ಹಾಗೆ} \kn{ಅನಿಸಿತು}.
\end{quote}

\begin{quote}
\emph{ā rātri praḷaya banda hāge anisitu.}
\end{quote}

\begin{quote}
That night it felt as if the end of the world had come.
\end{quote}

In the old stories, the flood that ends the world; in talk, any huge disaster.

\begin{cousinweb}
\textbf{\kn{ಪ್ರಳಯ}} (\emph{praḷaya}) is borrowed from Sanskrit \emph{pralaya}, \textquotedblleft{}dissolution, destruction of the world\textquotedblright{} (CDIAL 8755).
\end{cousinweb}

\begin{grammarlens}[title={What you've built}]
One more word for telling what happened.
\end{grammarlens}

\subsection*{Your turn}
\begin{itemize}
\raggedright
  \item \emph{Say it:} \emph{praḷaya}
  \item \emph{Say it:} \emph{praḷaya}, once more
  \item \emph{Say it:} \emph{ā rātri praḷaya banda hāge anisitu.} — That night it felt as if the end of the world had come.
  \item \emph{Recall:} say the Kannada for to reel, then the Kannada for to tremble, then say \emph{praḷaya} again
\end{itemize}

\begin{tcolorbox}[breakable,colback=teal!4,colframe=teal!35!black,title={Before you move on}]
What does \kn{ಪ್ರಳಯ} mean? (\textquotedblleft{}A deluge; a catastrophe\textquotedblright{}.) Say it once more.
\end{tcolorbox}
\section[\texorpdfstring{\kn{ಗಾಯಾಳು} (gāyāḷu)}{ಗಾಯಾಳು (gāyāḷu)}]{\kn{ಗಾಯಾಳು} (gāyāḷu) — an injured person; a casualty}

\label{lesson:KA-C346-gayalu}



\begin{quote}
Before the new one: say the Kannada for to tremble, then the Kannada for a deluge.
\end{quote}

\subsection*{You'll want to know: \kn{ಗಾಯಾಳು}}
\textbf{\kn{ಗಾಯಾಳು}} — \emph{gāyāḷu} — \textquotedblleft{}an injured person; a casualty\textquotedblright{}.

\begin{quote}
\kn{ಸ್ಫೋಟದ} \kn{ನಂತರ} \kn{ಗಾಯಾಳುಗಳನ್ನು} \kn{ಆಸ್ಪತ್ರೆಗೆ} \kn{ಕರೆದುಕೊಂಡು} \kn{ಹೋದರು}.
\end{quote}

\begin{quote}
\emph{sphōṭada nantara gāyāḷugaḷannu āspatrege karedukoṇḍu hōdaru.}
\end{quote}

\begin{quote}
After the explosion the injured were taken to hospital.
\end{quote}

Plural: \kn{ಗಾಯಾಳುಗಳು}, \emph{gāyāḷugaḷu}.

\kn{ಗಾಯ}, a wound, and \kn{ಆಳು}, a person: one who is hurt. A news word after accidents and blasts.

\begin{cousinweb}
\textbf{\kn{ಗಾಯಾಳು}} (\emph{gāyāḷu}) is a Kannada compound of \kn{ಗಾಯ}, \textquotedblleft{}a wound\textquotedblright{}, which goes back to Sanskrit \emph{ghāta}, \textquotedblleft{}a blow, a bruise\textquotedblright{} (CDIAL 4460), and the native Dravidian \emph{āḷ}, \textquotedblleft{}a person\textquotedblright{} (DEDR 399): a wounded person.
\end{cousinweb}

\begin{grammarlens}[title={What you've built}]
One more word for telling what happened.
\end{grammarlens}

\subsection*{Your turn}
\begin{itemize}
\raggedright
  \item \emph{Say it:} \emph{gāyāḷu}
  \item \emph{Say it:} \emph{gāyāḷu}, once more
  \item \emph{Say it:} \emph{sphōṭada nantara gāyāḷugaḷannu āspatrege karedukoṇḍu hōdaru.} — After the explosion the injured were taken to hospital.
  \item \emph{Recall:} say the Kannada for to tremble, then the Kannada for a deluge, then say \emph{gāyāḷu} again
\end{itemize}

\begin{tcolorbox}[breakable,colback=teal!4,colframe=teal!35!black,title={Before you move on}]
What does \kn{ಗಾಯಾಳು} mean? (\textquotedblleft{}An injured person; a casualty\textquotedblright{}.) Say it once more.
\end{tcolorbox}
